\documentclass[a4paper,11pt]{ltjsarticle}
\usepackage[haranoaji]{luatexja-preset}
\usepackage{amsmath,amssymb}
\usepackage[top=12mm,bottom=10mm,left=14mm,right=14mm]{geometry}
\usepackage{array}
\usepackage{tikz}
\usetikzlibrary{calc}
\pagestyle{empty}
\setlength{\parindent}{0pt}
\newlength{\qcol}\setlength{\qcol}{0.47\textwidth}
\newlength{\qgap}
\newlength{\anscolw}\setlength{\anscolw}{0.30\textwidth}
\newlength{\ansh}
\newcommand{\Q}[2]{\parbox[t]{\qcol}{\raggedright (#1)\hspace{0.6em}$\displaystyle #2$}}
\newcommand{\QR}[4]{\Q{#1}{#2}\hfill\Q{#3}{#4}\par\vspace{\qgap}}
\newcommand{\anscell}[1]{\rule[-\ansdrop]{0pt}{\ansh}(#1)}
\newlength{\ansdrop}

\begin{document}
{\large\bfseries 数学テスト　18}\hfill 名前(\underline{\hspace{32mm}})\par\vspace{3mm}
■（1）〜（12）の計算をしなさい。（13）、（14）は連立方程式を解きなさい。\par\vspace{3mm}
\setlength{\qgap}{11mm}
\QR{1}{(5x+2y)+(-3x-6y)}{2}{(a^2+2a-1)-(-a^2+4a+2)}
\QR{3}{36a^2b^3\div4a^2b^2\times3a}{4}{(12x+18y)\div6}
\QR{5}{(-3x)^2\times(-2x)}{6}{20a^2b^2\div(-5ab)}
\QR{7}{\dfrac{3a-b}{4}-\dfrac{a-2b}{2}+\dfrac{2a-b}{3}}{8}{10a-5-\{7b+(4a-3b)+2\}}
\QR{9}{3(2x-3y)+2(3x-2y)}{10}{(4x+y)-(3x-y)-(2x-y)}
\QR{11}{\dfrac{5}{3}(9x-12y)}{12}{45ab^3\div3b\div5ab}
\QR{13}{\left\{\begin{array}{l} 3x-2y=5 \\ 2x+3y=-1 \end{array}\right.}{14}{\left\{\begin{array}{l} x+\dfrac{y}{3}=1 \\ x+y=7 \end{array}\right.}
\vfill
\setlength{\ansh}{12mm}\setlength{\ansdrop}{8mm}%
\begin{center}\begin{tabular}{|p{\anscolw}|p{\anscolw}|p{\anscolw}|}\hline
\anscell{1} & \anscell{2} & \anscell{3} \\\hline
\anscell{4} & \anscell{5} & \anscell{6} \\\hline
\anscell{7} & \anscell{8} & \anscell{9} \\\hline
\anscell{10} & \anscell{11} & \anscell{12} \\\hline
\anscell{13} & \anscell{14} &  \\\hline
\end{tabular}\end{center}

\end{document}
